\documentclass[a4paper, serif, 10pt]{moderncv}

%% moderncv
% style and color
\moderncvstyle{banking}
\moderncvcolor{black}

% renew moderncv command to adapt for CJK colon
\renewcommand*{\cvitem}[3][.25em]{%
  \ifstrempty{#2}{}{\hintstyle{#2}：}{#3}%
  \par\addvspace{#1}}

\renewcommand*{\cvitemwithcomment}[4][.25em]{%
  \savebox{\cvitemwithcommentmainbox}{\ifstrempty{#2}{}{\hintstyle{#2}：}#3}%
  \setlength{\cvitemwithcommentmainlength}{\widthof{\usebox{\cvitemwithcommentmainbox}}}%
  \setlength{\cvitemwithcommentcommentlength}{\maincolumnwidth-\separatorcolumnwidth-\cvitemwithcommentmainlength}%
  \begin{minipage}[t]{\cvitemwithcommentmainlength}\usebox{\cvitemwithcommentmainbox}\end{minipage}%
  \hfill% fill of \separatorcolumnwidth
  \begin{minipage}[t]{\cvitemwithcommentcommentlength}\raggedleft\small\itshape#4\end{minipage}%
  \par\addvspace{#1}}

%% page layout/margins
\usepackage[top=2.5cm, bottom=2.5cm, left=1.5cm, right=1.5cm]{geometry}

\nopagenumbers{}

%% Babel config for English language
\usepackage[english]{babel}

%% fontspec
\usepackage{fontspec}

\IfFontExistsTF{Linux Libertine}{
  \setmainfont[Ligatures={TeX, Common}, Numbers=Lining, ItalicFont=Linux Libertine]{Linux Libertine}
}{}
\IfFontExistsTF{Linux Libertine O}{
  \setmainfont[Ligatures={TeX, Common}, Numbers=Lining, ItalicFont=Linux Libertine O]{Linux Libertine O}
}{}

%% CTeX
% CJK support, used to show CJK characters in the resume
%
% - fontset=none: disable builtin fontset but instead set the CJK font manually
% - heading=false: disable ctex heading
% - punct=kaiming: use kaiming punctuations styles for CJK
% - scheme=plain: use plain scheme, do not override \normalsize font size
% - space=auto: space settings for CJK characters
%
% ref:
% - http://ctan.mirrorcatalogs.com/language/chinese/ctex/ctex.pdf
\usepackage[UTF8, heading=false, punct=kaiming, scheme=plain, space=auto]{ctex}

\IfFontExistsTF{Noto Serif CJK SC}{
  \setCJKmainfont{Noto Serif CJK SC}
}{}
\IfFontExistsTF{Noto Sans CJK SC}{
  \setCJKsansfont{Noto Sans CJK SC}
}{}

%% line spacing
\usepackage{setspace}
\setstretch{1.125}

%% URL styling - use normal text instead of monospace
\urlstyle{same}

% Auto-underline all links
% Defer until \begin{document} so hyperref is already loaded
\AtBeginDocument{%
  \let\oldhref\href
  \renewcommand{\href}[2]{\oldhref{#1}{\underline{#2}}}%
}

%% Patch \httplink and \httpslink to handle full URLs
% The original moderncv macros always prepend http:// or https:// to URLs.
% This causes issues when the URL already has a protocol (e.g., https://example.com)
% resulting in malformed URLs like https://https://example.com.
% This patch checks if the URL already contains "://" and uses it directly if so.
% Note: We use \str_set:Nx (x-expansion) to expand macros like \@homepage first.
\ExplSyntaxOn
\str_new:N \g_yamlresume_url_str

\RenewDocumentCommand{\httpslink}{O{}m}{
  \str_set:Nx \g_yamlresume_url_str {#2}
  \str_if_in:NnTF \g_yamlresume_url_str {://}
    { \url{#2} }
    { \url{https://#2} }
}

\RenewDocumentCommand{\httplink}{O{}m}{
  \str_set:Nx \g_yamlresume_url_str {#2}
  \str_if_in:NnTF \g_yamlresume_url_str {://}
    { \url{#2} }
    { \url{http://#2} }
}
\ExplSyntaxOff

\name{林怡君}{}
\title{資深使用者經驗設計師 ｜ 打造以人為本的數位體驗}
\phone[mobile]{+886 912 345 678}
\email{yichun.lin@example.com}
\homepage{https://www.behance.net/yichunlin}

\address{中國臺灣，大安區，台北市，信義路二段 100 號 8 樓，106}{}{}

\extrainfo{{\small \faBehance}\ \href{https://www.behance.net/yichunlin}{@yichunlin} {} {} {} • {} {} {}
{\small \faLinkedin}\ \href{https://www.linkedin.com/in/yichunlin}{@yichunlin} {} {} {} • {} {} {}
{\small \faMedium}\ \href{https://medium.com/@yichunlin}{@yichunlin}}

\begin{document}

\maketitle

\section{簡介}

\cvline{}{擁有 8 年以上產品設計與使用者研究經驗，專精於 Web 與 App 的 UX 策略、互動設計與設計系統。
%
\begin{itemize}
\item 主導超過 20 個跨國與本土專案，涵蓋金融科技、電商、SaaS 與醫療照護領域
\item 擅長以質化與量化研究發掘使用者需求，並將研究洞見轉化為具體設計決策
\item 熟悉 Design Sprint、雙鑽石流程與敏捷開發，能有效與產品、工程團隊協作
\item 透過設計系統與反覆原型測試，協助團隊提升開發效率與產品易用性
\end{itemize}}
\section{教育背景}

\cventry{2012年9月–2014年7月}
        {碩士，設計系 互動設計組，成績：3.9}
        {國立臺灣科技大學}
        {\url{https://www.ntust.edu.tw}}
        {}
        {\begin{itemize}
\item 碩士論文以「高齡者行動支付應用程式之介面設計研究」獲選為系所優秀論文
\item 協助教授執行科技部研究計畫，負責使用者訪談、數據整理與原型測試
\end{itemize}
\textbf{課程}：人機互動、認知心理學、使用者經驗研究、互動多媒體設計、服務設計}

\cventry{2008年9月–2012年7月}
        {學士，工業設計系，成績：3.7}
        {國立臺北科技大學}
        {\url{https://www.ntut.edu.tw}}
        {}
        {\begin{itemize}
\item 畢業專題「城市公共座椅設計」入圍全國設計專題競賽前 10 名
\item 在校期間參與跨系所合作案，學習將工程技術與設計思維整合
\end{itemize}
\textbf{課程}：產品設計、人因工程、色彩學、設計素描、材料與製造程序}
\section{工作經歷}

\cventry{2021年1月至今}
        {資深 UX Designer}
        {澄心設計顧問股份有限公司}
        {\url{https://www.chengxin-design.example.com}}
        {}
        {\begin{itemize}
\item 主導金融科技與醫療照護產品的 UX 策略，包含使用者研究、資訊架構、互動流程與設計系統建置
\item 與產品經理、前後端工程師協作，透過設計衝刺與快速原型測試，將新功能上線時間縮短 30\%
\item 建立以使用者為中心的設計流程，導入易用性測試與數據分析，提升產品滿意度與留存率
\item 擔任設計導師，協助資淺設計師提升研究與視覺傳達能力
\end{itemize}
\textbf{關鍵字}：使用者研究、設計策略、設計系統、跨職能協作、敏捷開發}

\cventry{2018年7月–2020年12月}
        {UX Designer}
        {FutureLab 科技股份有限公司}
        {\url{https://www.futurelab.example.com}}
        {}
        {\begin{itemize}
\item 負責 B2B SaaS 產品線的介面設計與使用流程優化，涵蓋資料儀表板、權限管理與報表系統
\item 執行超過 15 場使用者訪談與遠端易用性測試，持續將研究發現轉化為設計迭代
\item 與工程團隊共同建立元件庫與設計規範，降低跨團隊溝通成本並提升開發效率
\item 參與產品需求定義與優先級排序，提出以數據驅動的設計方案，協助達成年度留存目標
\end{itemize}
\textbf{關鍵字}：B2B SaaS、易用性測試、設計系統、數據分析、使用者訪談}

\cventry{2016年8月–2018年6月}
        {UI/UX Designer}
        {好點子數位設計有限公司}
        {\url{https://www.goodidea.example.com}}
        {}
        {\begin{itemize}
\item 負責電商品牌與行銷活動網頁的視覺設計、RWD 切版與 Prototype 製作
\item 與設計總監合作建立設計規範，統一元件樣式與互動行為，提升品牌一致性
\item 運用 Google Analytics 與 Hotjar 分析使用者行為，提出頁面優化建議並提升轉換率 15\%
\end{itemize}
\textbf{關鍵字}：響應式設計、視覺設計、Prototype、電商、數據分析}
\section{語言}

\cvline{中文}{母語或雙語水平 \hfill \textbf{關鍵字}：母語、繁體中文}
\cvline{英語}{完全專業水平 \hfill \textbf{關鍵字}：TOEIC 900、產品簡報、跨國協作}
\section{專業技能}

\cvline{使用者研究}{專家 \hfill \textbf{關鍵字}：深度訪談、脈絡探查、易用性測試、問卷調查、Persona、顧客旅程地圖}
\cvline{互動設計}{專家 \hfill \textbf{關鍵字}：線框圖、互動原型、資訊架構、使用者流程、微互動}
\cvline{視覺設計}{高級 \hfill \textbf{關鍵字}：Figma、Adobe XD、Sketch、設計系統、排版與色彩}
\cvline{原型開發}{中級 \hfill \textbf{關鍵字}：HTML/CSS、JavaScript、Framer、Principle、ProtoPie}
\section{榮譽}

\cventry{2022年11月}
        {台灣設計研究院}
        {金點設計獎}
        {}
        {}
        {以「記憶守護者」失智症照護 App 榮獲年度最佳設計獎，肯定其在介面設計、服務流程與社會關懷上的創新。}

\cventry{2014年7月}
        {國立臺灣科技大學}
        {設計系所優秀論文獎}
        {}
        {}
        {碩士論文「高齡者行動支付應用程式之介面設計研究」獲選為當年度優秀論文。}
\section{證書}

\cventry{2023年1月}
        {Google}
        {Google UX Design Professional Certificate}
        {\url{https://www.coursera.org/}}
        {}
        {}

\cventry{2022年3月}
        {Nielsen Norman Group}
        {NN/g UX Certification}
        {\url{https://www.nngroup.com/}}
        {}
        {}
\section{出版刊物}

\cventry{2023年7月}
        {以服務設計觀點探討智慧健康照護平台之使用者經驗}
        {設計學報}
        {\url{https://www.jodesign.org.tw}}
        {}
        {探討智慧健康照護平台在使用者經驗上的挑戰，提出以服務藍圖與使用者旅程地圖整合的設計評估架構，並以實際案例驗證其可行性。}
\section{項目}

\cventry{2021年3月–2022年12月}
        {為失智症患者與家屬設計的照護協作 App，整合用藥提醒、日常紀錄與緊急聯絡功能。}
        {記憶守護者（Memory Keeper）}
        {\url{https://www.memorykeeper.example.com}}
        {}
        {\begin{itemize}
\item 主導田野調查與訪談，深入理解照顧者與患者的痛點，建立人物誌與使用情境
\item 設計完整的服務流程與互動原型，並進行 3 輪易用性測試，任務完成率提升 42\%
\item 與工程團隊密切協作，確保設計系統在 iOS 與 Android 平台上的一致體驗
\end{itemize}
\textbf{關鍵字}：服務設計、使用者研究、原型測試、失智症照護、設計系統}

\cventry{2020年1月–2021年6月}
        {重新設計醫療院所遠距掛號與視訊看診流程，降低患者的操作焦慮並提升預約成功率。}
        {遠距醫療預約平台 Redesign}
        {\url{https://www.telehealth.example.com}}
        {}
        {\begin{itemize}
\item 分析既有使用者流程與客服數據，找出 10 個關鍵斷點並進行重新設計
\item 建立符合無障礙規範的設計元件，確保高齡使用者也能順利操作
\item 系統上線後，預約成功率提升 27\%，客服進線量減少 18\%
\end{itemize}
\textbf{關鍵字}：醫療設計、無障礙設計、使用者流程、數據分析}
\section{興趣愛好}

\cvline{設計寫作}{Medium、設計方法、個案研究}
\cvline{戶外運動}{登山、單車、露營}
\section{志願服務}

\cventry{2019年1月–2022年12月}
        {實體活動志工}
        {台灣使用者經驗設計協會}
        {\url{https://uxpataiwan.example.com}}
        {}
        {\begin{itemize}
\item 協助籌辦年度使用者經驗研討會，負責講者邀請、工作坊規劃與現場活動執行
\item 參與線上社群文章翻譯與推廣，促進台灣設計社群知識交流
\end{itemize}}

\end{document}